\subsection{Gaussian Elimination}
Gaussian elimination is a method for solving systems of linear equations by transforming the system's augmented matrix into a simpler form using elementary row operations.
The goal is to reach \textbf{row echelon form} or \textbf{reduced row echelon form}, from which the solutions can be easily obtained.

\subsubsection*{Row Echelon Form} 
\paragraph*{Definition} A matrix is in \textbf{row echelon form} if it satisfies the following conditions:
\begin{itemize}
    \item All non-zero rows are above any rows of all zeros.
    \item The leading entry (the first non-zero number from the left, also called a \textbf{pivot}) of each non-zero row is to the right of the leading entry of the previous row.
    \item The leading entry in any non-zero row is 1 (this condition is sometimes omitted in basic row echelon form).
\end{itemize}

\paragraph*{Example} of a matrix in row echelon form:
\[\begin{bmatrix}
1 & 2 & -1 & | & 3 \\
0 & 1 & 3 & | & 4 \\
0 & 0 & 1 & | & 2 \\
\end{bmatrix}\]

It is often useful to further refine this to \textbf{reduced row echelon form}, which additionally requires that
the leading 1 in each non-zero row is the only non-zero entry in its \textbf{column}.

\paragraph*{Example} of a matrix in reduced row echelon form:
\[\begin{bmatrix}
1 & -2 & 0 & | & 2 \\
0 & 0 & 1 & | & -1 \\
0 & 0 & 0 & | & 0 \\
\end{bmatrix}\]

\paragraph*{Notes}
\begin{itemize}
    \item Every matrix can be reduced to row echelon form and reduced row echelon form using a sequence of elementary row operations.
    \item The reduced row echelon form of any matrix is unique.
    \item Row echelon form is not unique; different sequences of row operations can lead to different row echelon forms for the same matrix.
\end{itemize}

\paragraph*{Zero Matrix} A \textbf{zero matrix} is a matrix in which all entries are zero. It is often denoted by $0$. The solution
to the zero matrix is $\mathbb{R}^n$, where n is the number of variables.

\subsubsection*{Matrix Reduction Algorithm}
\begin{enumerate}
    \item Identify the leftmost non-zero column (the pivot column).
    \item Select a non-zero entry in the pivot column as the pivot. If necessary, swap rows to move this entry to the top of the pivot column.
    \item Scale the pivot row so that the pivot entry becomes 1.
    \item Use row addition to create zeros in all positions below the pivot.
    \item Cover the row containing the pivot and repeat the process for the submatrix that remains until all rows have been processed.
    \item (Optional) To achieve reduced row echelon form, repeat the process to create zeros above each pivot as well.
\end{enumerate}


\subsubsection*{Gaussian Elimination Algorithm}
\begin{enumerate}
    \item Augemnted Matrix -> RREF with Row Operations
    \item If a row of the form [0 0 ... 0 | b] with b non-zero exists, then the system is inconsistent (no solutions).
    \item Otherwise, assign free variables (if any) and solve for leading variables to find the solution set.
    \item Write leading variables in terms of the parameters
    \item If there are no nonleading variables, we gave 1 solution. If there are nonleading variables, we have infinitely many solutions.
\end{enumerate}



\paragraph*{Example 1} Solve the system:
\begin{align*}
    \begin{cases}
        3x + y - 4z = 1 \\
        x + 10z = 5 \\
        4x + y + 6z = 1 \\
    \end{cases} = &\begin{bmatrix}
        3 & 1 & -4 & | & 1 \\
        1 & 0 & 10 & | & 5 \\
        4 & 1 & 6 & | & 1 \\
    \end{bmatrix} \xmapsto{R2 \leftrightarrow R1} \begin{bmatrix}
        1 & 0 & 10 & | & 5 \\
        3 & 1 & -4 & | & 1 \\
        4 & 1 & 6 & | & 1 \\
    \end{bmatrix} \xmapsto{R2 - 3R1 \mapsto R2} \begin{bmatrix}
        1 & 0 & 10 & | & 5 \\
        0 & 1 & -34 & | & -14 \\
        4 & 1 & 6 & | & 1 \\
    \end{bmatrix}\\ 
    \xmapsto{R3 - 4R1 \mapsto R3} &\begin{bmatrix}
        1 & 0 & 10 & | & 5 \\
        0 & 1 & -34 & | & -14 \\
        0 & 1 & -34 & | & -19 \\
    \end{bmatrix} 
    \xmapsto{R3 - R2 \mapsto R3} \begin{bmatrix}
       1 & 0 & 10 & | & 5 \\
        0 & 1 & -34 & | & -14 \\
        0 & 0 & 0 & | & -5 \\
    \end{bmatrix} \implies \text{ no solutions}
\end{align*}

\paragraph*{Example 2} Solve the system:
\begin{align*}
    \begin{cases}
        x_1 - 2x_2 - x_3 + 3x+4 = 1 \\
        2x_1 - 4x_2 + x_3 = 5 \\
        x_1 - 2x_2 + 2x_3 -3x_4 = 4 \\
    \end{cases} = & \begin{bmatrix}
        1 & -2 & -1 & 3 & | & 1 \\
        2 & -4 & 1 & 0 & | & 5 \\
        1 & -2 & 2 & -3 & | & 4 \\
    \end{bmatrix} \xmapsto{R2 - 2R1 \mapsto R2} \begin{bmatrix}
        1 & -2 & -1 & 3 & | & 1 \\
        0 & 0 & 3 & -6 & | & 3 \\
        1 & -2 & 2 & -3 & | & 4 \\
    \end{bmatrix} \\
    \xmapsto{R3 - R1 \mapsto R3} & \begin{bmatrix}
        1 & -2 & -1 & 3 & | & 1 \\
        0 & 0 & 3 & -6 & | & 3 \\
        0 & 0 & 3 & -6 & | & 3 \\
    \end{bmatrix} \xmapsto{R3 - R2 \mapsto R3} \begin{bmatrix}
        1 & -2 & -1 & 3 & | & 1 \\
        0 & 0 & 3 & -6 & | & 3 \\
        0 & 0 & 0 & 0 & | & 0 \\
    \end{bmatrix} \\
    \xmapsto{\frac{1}{3}R2 \mapsto R2} &\begin{bmatrix}
        1 & -2 & -1 & 3 & | & 1 \\
        0 & 0 & 1 & -2 & | & 1 \\
        0 & 0 & 0 & 0 & | & 0 \\
    \end{bmatrix} = \begin{cases}
        x_1 - 2x_2 - x_3 + 3x_4 = 1 \\
        x_3 - 2x_4 = 1 \\
    \end{cases} \\
    \implies & \begin{cases}
        x_2 = s \\
        x_4 = t \\
        x_1 = 1 + 2s + t \\
        x_3 = 1 + 2t \\
    \end{cases}
\end{align*}

\hrulefill 

\subsubsection*{Rank of a Matrix}
\paragraph*{Definition} The \textbf{rank} of a matrix is the number of leading 1's (pivots) in its row echelon form or reduced row echelon form.
It represents the dimension of the column space (or row space) of the matrix, which is the span of its columns (or rows). Rank is useful 
for finding the number of parameters in the solution set of a linear system. 

\paragraph*{Example 3} Consider the matrix from Example 2. Find the rank and number of parameters:
\begin{align*}
    A &= \begin{bmatrix}
        1 & -2 & -1 & 3 & | & 1 \\
        2 & -4 & 1 & 0 & | & 5 \\
        1 & -2 & 2 & -3 & | & 4 \\
    \end{bmatrix} \xmapsto{\text{RREF(A)}} \begin{bmatrix}
        1 & -2 & 0 & 1 & | & 0 \\
        0 & 0 & 1 & -2 & | & 1 \\
        0 & 0 & 0 & 0 & | & 0 \\
    \end{bmatrix} \implies \text{rank}(A) = 2 \\
    n &= 4 \text{ (number of variables) } \implies \text{ number of parameters } = n - \text{rank}(A) = 4 - 2 = 2
\end{align*}

\paragraph*{Example 4} Find the rank of the following matricies:
\begin{align*}
    A &= \begin{bmatrix}
        1 & 0 & 0 & 2 & 0 & 0 \\
        0 & 0 & 0 & 1 & 3 & 0 \\
        0 & 0 & 0 & 0 & 0 & 0 \\
    \end{bmatrix} \implies \text{rank}(A) = 2 \\
    B &= \begin{bmatrix}
        0 & 0 \\
        0 & 3 \\
    \end{bmatrix} \xmapsto{R1 \leftrightarrow R2} \begin{bmatrix}
        0 & 3 \\
        0 & 0 \\
    \end{bmatrix} \xmapsto{\frac{1}{3}R1 \mapsto R1} \begin{bmatrix}
        0 & 1 \\
        0 & 0 \\
    \end{bmatrix} \implies \text{rank}(B) = 1 \\
    C &= \begin{bmatrix}
        0 & 0 & 0 \\
        0 & 0 & 0 \\
    \end{bmatrix} \implies \text{rank}(C) = 0
\end{align*}